\subsection{Auxiliary squares and a parity exchange}

We next give a sufficient condition on the owner graph when all auxiliary
components are four-cycles. The condition concerns the owner dissociation
number, which has an explicit formula for paths and even cycles.

\begin{lemma}\label{lem:auxiliary-square-normalization}
Suppose $A=P\cup Q$ is a disjoint union of four-cycles on the index set $V$
of a partial template, with $|V|=2b$. Then
\[
 \beta_I(H)=2b+
 \max\{|S|:A[S]\text{ is edgeless and }B[S]\text{ is dissociation}\}.
\]
\end{lemma}
\begin{proof}
For a permitted deleted set $U$ in Theorem~\ref{thm:all-helper}, put
$r_K=\nu(A[U\cap K])$ on each auxiliary square $K$.
Its owner complement $K\setminus U$ contains an independent set of size
$2-r_K$. If $r_K=0$, extend $U\cap K$ to an opposite pair and take the other
pair. If $r_K=1$, the complement is nonempty and one vertex suffices.
If $r_K=2$, the complement is empty. Their union $S$ is independent in $A$
and lies in the dissociation set $V\setminus U$ of $B$, so it remains
dissociation in $B$. Since there are $b/2$ squares,
$|S|=b-\nu(A[U])$.

Conversely, an $A$-independent set $S$ meets each square in zero, one or two
opposite vertices. For $U=V\setminus S$ the matching cost there is exactly
$2-|S\cap K|$. If $B[S]$ is dissociation this $U$ is permitted, and its total
cost is $b-|S|$. Taking the two optima and substituting in
Theorem~\ref{thm:all-helper} proves the identity.
\end{proof}

The next statement is purely combinatorial; it needs no compatibility with
the physical bipartition of a template.

\begin{theorem}\label{thm:auxiliary-square-parity}
Let $A$ be a disjoint union of four-cycles on a set $V$ of size $N=4s$.
Let $B$ be a finite simple graph on $V$ whose components are paths or even
cycles. If $\diss(B)\ge2\lceil N/3\rceil-2$, then there is an $A$-independent
set $S$ with $B[S]$ dissociation and $|S|\ge\lceil N/3\rceil$.
Consequently every physical partial template with these auxiliary and owner
graphs satisfies $\beta_I(H)\ge\lceil8b/3\rceil$.
\end{theorem}
\begin{proof}
For a path or a cycle of order $m$,
\[
 \diss(P_m)=\lceil2m/3\rceil,\qquad
 \diss(C_m)=\lfloor2m/3\rfloor.
\]
For paths, disjoint consecutive triples require $\lfloor m/3\rfloor$
omissions, attained by deleting positions $3,6,\ldots$.
For cycles, every omission meets at most three cyclic triples, so at least
$\lceil m/3\rceil$ omissions are required. Cyclic gaps of size at most three
between that many omissions attain the bound. These choices add over the
components of $B$.

We first record an exchange on an even cycle of length $m$ not divisible by
three. For any prescribed consecutive vertices $x,y$, there are two maximum
dissociation sets differing exactly on $\{x,y\}$.
For $m=4$, any two vertices form a maximum set, so the assertion is immediate.
For $m>4$, put $k=\lceil m/3\rceil$ and $\delta=3k-m\in\{1,2\}$.
Choose a deleted set of size $k$ whose cyclic gaps immediately before and
after $x$ are $2,3$, respectively. All remaining gaps are three, except for
one additional gap two when $\delta=2$.
Their sum is $3k-\delta=m$; there is space for the additional gap because
$m>4$ gives $k\ge3$ when $\delta=1$, and $k\ge4$ when $\delta=2$.
Replacing $x$ by its successor $y$ changes the two local gaps from $2,3$ to
$3,2$. Both deleted sets still meet every consecutive triple and have minimum
size $k$. Taking complements proves the exchange.

Choose a maximum dissociation set $D$ of $B$. In each auxiliary square $K$,
\[
 \alpha(A[D\cap K])\ge\left\lceil\frac{|D\cap K|}{2}\right\rceil.
\]
Put $h=\lceil N/3\rceil$. If $|D|\ge2h-1$, these square inequalities
already give an independent subset of size at least $h$.
Otherwise the hypothesis forces $|D|=2h-2$, an even integer.
For every nonzero $N$ divisible by four,
\[
 2\lceil N/3\rceil-2<\lceil2N/3\rceil;
\]
this follows immediately in the residues $N=6t,6t+2,6t+4$.
There is therefore an owner cycle whose length is not divisible by six.
Otherwise every owner component has dissociation number at least two thirds
of its order, giving $|D|\ge\lceil2N/3\rceil$, a contradiction.
Choose such a cycle $C$.

If $C$ spans two auxiliary squares, it has an edge $xy$ joining different
squares. Apply the prescribed-edge exchange on $C$, keeping maximum choices
in every other owner component. This gives two global maximum dissociation
sets $D_0,D_1$ with symmetric difference $\{x,y\}$.
If $D_0$ meets every square evenly, $D_1$ meets the two squares of $x,y$
oddly. Otherwise $D_0$ already has at least two odd intersections, because
its total order is even. For the resulting set $D_*$, the square inequalities
therefore give
\[
 \alpha(A[D_*])\ge |D_*|/2+1=h.
\]
The intersections and the two rounding gains belong to distinct auxiliary
components and are counted once.

If $C$ lies in one auxiliary square, it has order four and fills that square.
Choose an auxiliary opposite pair as its maximum owner dissociation set,
and retain maximum choices in all other owner components. That square
contributes two independent vertices instead of the half-cardinality bound
one. This again gives $\alpha(A[D_*])\ge|D_*|/2+1$.

In either case choose $S$ independent in $A[D_*]$. It remains dissociation
in $B$ by restriction. Lemma~\ref{lem:auxiliary-square-normalization}, with
$N=2b$, gives the template conclusion. The empty case is immediate.
\end{proof}

The owner condition can be checked from
\[
 \diss(B)=\sum_{P_m}\lceil2m/3\rceil+
                  \sum_{C_m}\lfloor2m/3\rfloor.
\]
For example, it permits arbitrary paths and $6k$-cycles together with one
$6k+4$-cycle, or with at most two $6k+2$-cycles, when their orders sum to $N$.
The total deficit relative to $2N/3$ is at most $2/3$, which gives
$\diss(B)\ge\lfloor2N/3\rfloor\ge2\lceil N/3\rceil-2$ by integrality. Simple even cycles here have order at least four; cycles of the
form $6k+2$ therefore start at order eight.

For a fixed witness extraction with these auxiliary squares and this owner
condition, the separate retained set combines with the template as before.
It gives $Z=3\diss(G)-3L-8b'\ge0$, hence $51(|S|-\diss(G))\le5|V(G)|$ by
the existing ledger. To bound $\SF-\diss$, the condition must hold for an
extraction of a maximum-cardinality minimal forming set.
The owner-size criterion is sufficient, and may fail even for templates
covered by the earlier square-owner theorem.
